\renewcommand{\sectionauthor}{Gustavo Ale}

\begin{multicols*}{2}

    \section*{Geometria}
    A geometria é provavelmente a área mais antiga da matemática, precursora da própria álgebra e sendo um dos pilares
    da matemática. A palavra geometria é resultado da combinação das palavras gregas \textit{geo} e
    \textit{metron}/\textit{metri}, significando respectivamente Terra e medição, pois a área da geometria consiste da
    medição e entendimento das relações e propriedades  contidas nas fíguras geométricas: comprimento, distância, ângulo,
    área, volume e perímetro, as fíguras geométricas por sua vez são parte integrante da natureza e da Terra (\textit{geo}).

    \subsection*{Definições gerais}
    Antes do estudo da geometria devemos ter ciência das definições de conceitos presentes nesse ramo da matemática,
    dentro destes conceitos se enquadram as medidas, os ângulos, as fíguras geométricas e suas componentes.

    \subsection*{Fíguras geométricas e suas componentes}
    A primeira componente das fíguras geométricas e a mais simples delas é o ponto, o mesmo deve ser imaginado como um
    ponto infinitesimal, não possuindo comprimento, área ou perímetro. Quando um ponto é integrante de uma fígura
    geométrica o mais comum é chamarmos ele de \textit{vértice}

    \begin{figure}[H]
        \centering
        \includegraphics[width=0.7\columnwidth]{assets/image_placeholder.png}
        \caption{Vértice de um triângulo }
    \end{figure}

    A linha por sua vez é o elemento que liga dois pontos $P_1$ e $P_2$ em uma reta, assim como o ponto é idealmente
    infinitesimal a linha idealmente não possui espessura. Quando uma linha é parte integrante de uma fígura geométrica
    sua denominação é aresta.

    \begin{figure}[H]
        \centering
        \includegraphics[width=0.7\columnwidth]{assets/image_placeholder.png}
        \caption{Aresta de um quadrado}
    \end{figure}


    Quando se tem 3 ou mais pontos o resultado da conexão deles é tido como superfície, plano ou fígura geométrica
    dependendo da área da matemática e do contexto de estudo, porém vale ter em mente que também existem superfícies
    não bidimensionais. Uma excessão a isso são os círculos e as elipses que não possuem vértices. No contexto da
    geometria a ser estudado nesse capítulo a conexão de 3 ou mais vértices geram exclusivamente fíguras geométricas.

    \begin{figure}[H]
        \centering
        \includegraphics[width=0.7\columnwidth]{assets/image_placeholder.png}
        \caption{Um quadrilátero com as arestas realçadas em vermelho}
    \end{figure}

    No contexto de fíguras geométricas tridimenssionais, ou chamados sólidos, os sólidos podem ser fígurados por 4 ou
    mais pontos ou através manipulação de fíguras geométricas bidimensionais num espaço tridimenssional.

    Grandezas notórias das componentes citadas:
    \begin{itemize}
        \setlength\itemsep{1.15pt}
        \item Ponto, vértice: distância relativo a outro objeto.
        \item Reta: comprimento, distância e ângulo relativos a outro objeto.
        \item Superfície: área, perímetro.
        \item Sólido: volume.
    \end{itemize}

    \begin{figure}[H]
        \centering
        \includegraphics[width=0.7\columnwidth]{assets/image_placeholder.png}
        \caption{Tetraedro, o objeto tridimenssional com menos faces}
    \end{figure}

    \subsection*{Grandezas e unidades de medida}
    \subsubsection*{Distância e comprimento}
    Distância e comprimento medem o quão longe dois pontos estão entre si, e no caso de arestas essa medida é
    denotada como comprimento. A unidade de medida base utilizada para comprimento é o metro (símbolo ${m}$), além do
    metro existem seus múltiplos que são igualmente utilizados dependendo da distância/comprimento aferido.
    \begin{table}[H]
        \resizebox{\columnwidth}{!}{%
            \begin{tabular}{l|l|l}
                \rowcolor{gray!25}
                \textbf{Nome} & \textbf{Sigla} & \textbf{Equivalência} \\ \hline
                picometro     & $pm$           & $10^{-12}m$           \\ \rowcolor{gray!10}
                nanometro     & $nm$           & $10^{-9}m$            \\
                micrometro    & $\mu m$        & $10^{-6}m$            \\ \rowcolor{gray!10}
                milimetro     & $mm$           & $10^{-3}m$            \\
                centímetro    & $cm$           & $10^{-2}m$            \\ \rowcolor{gray!10}
                decímetro     & $dm$           & $10^{-1}m$            \\
                metro         & $m$            & $1m$                  \\ \rowcolor{gray!10}
                decâmetro     & $dam$          & $10m$                 \\
                hêctometro    & $hm$           & $100m$                \\ \rowcolor{gray!10}
                quilometro    & $km$           & $1000m$
            \end{tabular}%
        }
        %\caption{Unidades medida de comprimento e distância}
    \end{table}
    Outras medidas de comprimento são: \\
    \textbf{Perímetro:} comprimento do contorno de uma fígura geométrica. \\
    \textbf{Raio:} distância entre o centro de uma circunferência até seu contorno ou superfície.\\
    \textbf{Diâmetro:} comprimento de reta que passe pelo centro da circunferência e cujo seus pontos de
    início e fim estejam sobre a circunferência\\
    \textbf{Circunferência:} perímetro de uma circunferência.\\

    \begin{figure}[H]
        \centering
        \includegraphics[width=0.7\columnwidth]{assets/image_placeholder.png}
        \caption{Raio, diâmetro e circunferência de um círculo}
    \end{figure}

    \subsubsection*{Área}
    Área é a medida que expressa a quantidade de espaço bidimenssional ocupado por uma fígura geométrica,
    área de superfície é o equivalente da área para uma superfície ou face de um objeto tridimenssional. A unidade base
    para área é o metro quadrado (símbolo ${m^2}$). Seus múltiplos também acompanham o termo 'quadrado',
    e a equivalência é a mesma do metro, porém elevada ao quadrado. Ex.: $1m = 10^2cm$ e $1m^2 = (10^2cm)^2 = 10^4cm^2$,
    lembrar essa regra pode facilitar na hora da conversão de área.

    \begin{table}[H]
        \resizebox{\columnwidth}{!}{%
            \begin{tabular}{l|l|l}
                \rowcolor{gray!25}
                \textbf{Nome}       & \textbf{Sigla} & \textbf{Equivalência}       \\ \hline
                picometro quadrado  & $pm^2$         & $10^{-24}m^2$               \\ \rowcolor{gray!10}
                nanometro quadrado  & $nm^2$         & $10^{-18}m^2$               \\
                micrometro quadrado & $\mu m$        & $10^{-12}m^2$               \\ \rowcolor{gray!10}
                milimetro quadrado  & $mm^2$         & $10^{-6}m^2$                \\
                centímetro quadrado & $cm^2$         & $10^{-4}m^2$                \\ \rowcolor{gray!10}
                decímetro quadrado  & $dm^2$         & $10^{-2}m^2$                \\
                metro quadrado      & $m^2$          & $1m^2$                      \\ \rowcolor{gray!10}
                decâmetro quadrado  & $dam^2$        &                             \\ \rowcolor{gray!10}
                are                 & $are$          & \multirow{-2}{*}{$100m^2$}  \\
                hêctometro quadrado & $hm^2$         &                             \\
                hectare             & $ha$           & \multirow{-2}{*}{$10^4m^2$} \\ \rowcolor{gray!10}
                quilometro quadrado & $km^2$         & $10^6m^2$
            \end{tabular}%
        }
        %\caption{Unidades medida de área}
    \end{table}

    \subsubsection*{Volume}
    Assim como a área expressa a quantidade de espaço bidimenssional ocupado por uma fígura geométrica, o volume expressa
    a quantidade de espaço tridimenssional ocupado por um objeto. No caso do volume existe duas unidades de medida base,
    o metro cúbico e o litro, ambos são usados regularmente, na qual o litro é a unidade mais usada para volumes
    pequenos, enquanto o metro cúbico é usado para expressar volumes maiores.

    \begin{table}[H]
        \resizebox{\columnwidth}{!}{%
            \begin{tabular}{l|l|l}
                \rowcolor{gray!25}
                \textbf{Nome}     & \textbf{Sigla} & \textbf{Equivalência}          \\ \hline
                milimetro cúbico  & $mm^3$         & $10^{-9}m^3$                   \\ \rowcolor{gray!10}
                centímetro cúbico & $cm^3$         &                                \\ \rowcolor{gray!10}
                mililitro         & $ml$           & \multirow{-2}{*}{$10^{-6}m^3$}  \\
                decímetro cúbico  & $dm^3$         &                                \\ %\hline
                litro             & $l$            & \multirow{-2}{*}{$10^{-3}m^3$} \\ \rowcolor{gray!10}
                metro cúbico      & $m^3$          & $1m^2$                         \\ %\hline                      
            \end{tabular}%
        }
        %\caption{Unidades medida de volume}
    \end{table}

    \subsubsection*{Ângulo}
    
    % teste matemática em duas colunas
    % \color{black}
    % \begin{align}
    %     f(x) &= \frac{sen(x)}{cos(x)} \\[10pt]
    %     f(x) &= tan(x)
    % \end{align}
    % \color{mediumgray}

    %\pagebreak
    % \newpage
    % \section*{Ângulos}
    % \lipsum[1]

    % \subsection*{Simetrias de figuras planas}
    % \lipsum[1-2]    

\end{multicols*}